%======================================================================
% GMI-04 Zobecneni variet.pdf
%======================================================================
\section{Zobecnění variet}
\strana{GMI-04b-1}
\pozn{titulní strana -- bez dalšího obsahu}

%----------------------------------------------------------------------
\newpage
\strana{GMI-04b-2}
\subsection{Fibrované prostory}

\begin{definice}[fibrovaný prostor (bundle)]
$\mathbb{E}M$ je \kw{fibrovaný prostor} se stand.\ fibrem $\mathbb{E}$ $\equiv$
\begin{itemize}[nosep]
  \item[] $\mathbb{E}M$ je dif.\ varieta \hfill -- totální prostor, bundle
  \item[] $M$ je dif.\ varieta \hfill -- podkladový [base] prostor
  \item[] $\mathbb{E}$ prostor s nějakou str.\ \hfill -- standardní vlákno [fiber]
  \item[] $\pi:\mathbb{E}M\to M$ \quad projekce dávající fibraci
  \item[] $\mathbb{E}_xM = \pi^{-1}x$ \quad vlákno [fiber] nad $x\in M$
  \item[] $\mathbb{E}_xM\cong\mathbb{E}$ \quad izomorfismus vlákna se stand.\ fibrem \\
          izomorfismus respektuje strukturu $\mathbb{E}$ \\
          přesněji: lokální izom.\ \quad $\mathbb{E}U\cong U\times\mathbb{E}$
\end{itemize}
\end{definice}

\begin{definice}[řezy fibrovaného prostoru]
\[
  \phi: M\to\mathbb{E}M \qquad x\mapsto\phi(x)\in\mathbb{E}_xM
\]
tj.\ \quad $\pi\circ\phi = \mathrm{id}$
\end{definice}

tečný a kotečný bundle jsou vekt.\ fibr.\ prostory \\
\quad $\mathcal{T}M$ je prostor řezů $\mathbb{T}M$ \\
\quad $\mathcal{T}^*M$ je prostor řezů $\mathbb{T}^*M$

tečný a kotečný prostor je svázán s $M$ více než obecný fibrovaný prostor nad $M$

\begin{vysvetleni}[Opakování]
Tato strana opakuje definici fibrovaného prostoru z GMI-04a (viz tamní vysvětlivky). Podstatné je \emph{lokální} ztotožnění $\mathbb{E}U\cong U\times\mathbb{E}$,
které respektuje projekci; globálně může být bandl netriviální. Řezy bandlu jsou \uv{pole} -- vektorová pole jsou řezy $\mathbb{T}M$ (Krtouš, odd.~2.4).
\end{vysvetleni}

%----------------------------------------------------------------------
\newpage
\strana{GMI-04b-3}
\subsection{Zobecnění variety pro obecnou strukturu}

dif.\ variety lze zobecnit pro případ, kdy souřadnice nabývají hodnot v obecném prostoru,
resp.\ připouštějí se pouze specifické transf.\ v prostoru souřadnic

\begin{definice}[pseudogrupa transformací]
\kw{pseudogrupa} $\Gamma$ \kw{transformací} na souřadnicovém prostoru $S$ $\equiv$
\begin{itemize}[nosep]
  \item[] $S$ \quad topol.\ prostor \hfill (prostor hodnot souřadnic -- např.\ $\mathbb{R}^n$)
  \item[] $\Gamma$ \quad množina transf.\ $\underset{\text{(otevř.}\to\text{otevř.)}}{S\to S}$ \hfill (přípustné transf.\ souř.\ -- např.\ $C^\infty$)
  \item[--] $\alpha\in\Gamma$ je spojité
  \item[--] $\alpha\in\Gamma \Rightarrow \forall\, U\subset S$ otevř.\ \quad $\alpha|_U\in\Gamma$
  \item[--] $\alpha$ homeomorf.\ na $U = \bigcup_\mu U_\mu$, \quad $U_\mu\subset S$ otevřené \\
            $\alpha|_{U_\mu}\in\Gamma \Rightarrow \alpha\in\Gamma$
  \item[--] $U\subset S$ otevř.\ $\Rightarrow \mathrm{id}|_U\in\Gamma$
  \item[--] $\alpha\in\Gamma \Rightarrow \alpha^{-1}\in\Gamma$
  \item[--] $\alpha,\beta\in\Gamma \Rightarrow \alpha\circ\beta\in\Gamma$ \quad na příslušné oblasti
\end{itemize}
\end{definice}

\begin{vysvetleni}[Proč \uv{pseudo}grupa]
Přechodová zobrazení mezi mapami nejsou definována na celém $S$, ale jen mezi otevřenými podmnožinami (obrazy překryvů map). Proto nejde o grupu:
složení $\alpha\circ\beta$ má smysl jen tam, kde se obor hodnot $\beta$ potkává s definičním oborem $\alpha$. Axiomy zaručují, že $\Gamma$ je uzavřena
na restrikce, inverze a skládání a že příslušnost do $\Gamma$ je \emph{lokální} vlastnost (lze ji ověřit na pokrytí -- třetí bod).
Příklady: všechny $C^k$ difeomorfismy otevřených množin $\mathbb{R}^n$; $C^\infty$ difeomorfismy s kladným Jacobiánem; biholomorfní zobrazení v $\mathbb{C}^n$.
\end{vysvetleni}

\begin{definice}[varieta se strukturou $\Gamma$]
varieta \dots\ $M$ se strukturou $\Gamma$ $\equiv$ \\
stejná def.\ jako pro dif.\ varietu, jen $\mathbb{R}^n$, $C^\infty$ nahrazeno strukturou $S$, $\Gamma$
\end{definice}

\begin{priklad}[příklady]
\begin{itemize}[nosep]
  \item \kw{orientovaná dif.\ varieta} \\
        -- připouští se pouze transf.\ souř.\ s $\det\dfrac{\partial y^i}{\partial x^j} > 0$
  \item \kw{komplexní variety} \\
        -- $S = \mathbb{C}^n$, holomorfní transf.
\end{itemize}
\end{priklad}

\begin{vysvetleni}[$\Gamma$-varieta]
Varieta se strukturou $\Gamma$ je topologický prostor s atlasem map do $S$, jehož všechna přechodová zobrazení $x_{(\alpha)}\circ x_{(\beta)}^{-1}$ leží v $\Gamma$
(a atlas je maximální). Volbou $\Gamma$ volíme, které vlastnosti jsou \uv{nezávislé na souřadnicích}: menší $\Gamma$ = bohatší struktura.
Orientovaná varieta: přechody s $\det\partial y^i/\partial x^j>0$ zachovávají orientaci báze, takže orientace je dobře definovaná (viz Krtouš, odd.~5.6).
Komplexní varieta: holomorfnost funkce je pak pojem nezávislý na mapě. Další příklady: reálně analytické, afinní či plochá struktura (přechody afinní).
\end{vysvetleni}

%----------------------------------------------------------------------
\newpage
\strana{GMI-04b-4}
\subsection{Varieta v řeči prostoru funkcí}

dif.\ varietu $M$ lze charakterizovat prostorem hl.\ fcí \\
několik příbuzných přístupů
\begin{itemize}[nosep]
  \item \kw{strukturní sheaf} dif.\ var.
  \item \kw{locally ringed space}
  \item \kw{$C^*$ algebry} s lokální hladkostí
\end{itemize}

zadání \uv{prostoru} fcí umožní zrekonstruovat top.\ či dif.\ strukturu na varietě či dokonce varietu samotnou

změnou funkcionálních vlastností pr.\ fcí lze zavést abstraktní pojem variety jako podkladu pro příslušný pr.\ fcí \\
umožňuje např.\ zavést
\begin{itemize}[nosep]
  \item[] \kw{supervariety} \quad (supersymetrie)
  \item[] \kw{nekomutativní geometrii}
\end{itemize}

\begin{vysvetleni}[Geometrie z funkcí]
Myšlenka: místo bodů a map zadáme algebru funkcí a body z ní zrekonstruujeme. Např.\ pro kompaktní Hausdorffův $X$ odpovídají body $x$ maximálním ideálům
$\{f: f(x)=0\}$ algebry $C(X)$ (Gelfandova dualita) a pro hladkou varietu odpovídají body homomorfismům $\mathcal{F}M\to\mathbb{R}$, $f\mapsto f(x)$.
Vektory v bodě jsou derivace této algebry (Krtouš, odd.~2.2: vektor jako diferenciální operátor 1.\ řádu).
Když pak algebru nahradíme obecnější (s antikomutujícími Grassmannovými prvky, resp.\ nekomutativní), dostaneme \uv{prostor}, který body v obvyklém smyslu
mít nemusí -- supervarieta, resp.\ nekomutativní geometrie.
\end{vysvetleni}

%----------------------------------------------------------------------
\newpage
\strana{GMI-04b-5}
\subsection{Svazky a okruhované prostory}

opakování z algebry:
\begin{itemize}
  \item \kw{abelovská grupa} \\
        operace $+$, asociat., komutat., inverze, prvek $0$
  \item \kw{komutativní těleso} $\mathbb{K}$ (pole) \quad [field] \\
        \begin{tabular}{@{}ll}
          operace $+$ & abelovská gr. \\
          \phantom{operace} $\cdot$ & asociat., komutat., distribut., prvek $1$ \\
                                     & inverze mimo $0$
        \end{tabular} \\
        typicky $\mathbb{K} = \mathbb{R},\mathbb{C}$
  \item \kw{vekt.\ prostor} $V$ nad $\mathbb{K}$ \\
        operace $+$ \quad abel.\ gr. \\
        násobení skalárem z $\mathbb{K}$ \quad -- linearita (tj.\ asoc.\ a distrib.)
  \item \kw{algebra} $A$ nad $\mathbb{K}$ \\
        vekt.\ pr.\ nad $\mathbb{K}$ s operací $+$ \\
        násobení v $A$ konzist.\ s operacemi vekt.\ pr.\ (bi-linearita)
        \begin{itemize}[nosep]
          \item[-] asociativní algebry
          \item[-] komutativní alg.
          \item[-] Lieovy alg.
        \end{itemize}
  \item \kw{okruh} $R$ \\
        obdobné tělesu, ale není zaručena existence inverze \\
        Př: funkce na varietě
  \item \kw{modul} $V$ nad $R$ \quad ($R$-modul $V$) \\
        obdobné vekt.\ pr., ale nad okruhem \\
        Př: vekt.\ a tenz.\ pole na varietě, řezy vekt.\ bandlů
  \item \kw{algebra} $A$ \kw{nad} \oprava{\kw{okruhem}}{modulem} $R$ \\
        obdobné algebře, jen nad okruhem \\
        Př: algebra fcí na varietě \quad (okruh $R$ je vždy alg.\ nad $R$) \\
        \phantom{Př:} operátory $\mathcal{T}^1_1M$ na varietě \dots
\end{itemize}

\begin{vysvetleni}[Proč okruhy a moduly]
Hladké funkce $\mathcal{F}M$ tvoří komutativní okruh (s jedničkou), ale ne těleso: funkce s nulovým bodem nemá inverzi, ačkoli není nulová.
Proto vektorová pole $\mathcal{T}M$ nejsou vektorový prostor nad $\mathcal{F}M$, ale $\mathcal{F}M$-modul -- násobit je lze funkcemi, dělit ne.
Okruh, nad kterým je algebra, je zde $R=\mathcal{F}M$ (odtud oprava \uv{nad okruhem}). Tenzorová pole $\mathcal{T}^1_1M$ tvoří algebru nad $\mathcal{F}M$
s násobením $(AB)^a{}_b=A^a{}_cB^c{}_b$ (skládání lineárních operátorů v každém bodě).
\end{vysvetleni}

%----------------------------------------------------------------------
\newpage
\strana{GMI-04b-6}
\begin{definice}[modul $V$ nad okruhem $R$ ($R$-moduly)]
$a_1,\dots,a_k\in V$ jsou \kw{lin.\ závislé}, pokud
\[
  r_1a_1 + \dots + r_ka_k = 0 \qquad\text{pro } \oprava{r_j\in R\text{ ne všechna nulová}}{r_j\in R}
\]
\pozor{nezaručuje možnost vyjádření $a_i$ pomocí ostatních!}

$e_i$ tvoří \kw{bázi} modulu $V$ \quad $\forall a\in V\ \exists a^i\in R$ \quad $a = \sum_i a^ie_i$
\end{definice}

\begin{tabular}{@{}ll}
\kw{volný modul} & -- existuje báze \\
\kw{konečně generovaný modul} & -- existuje konečná báze \\
\kw{duální modul} $V^*$ & -- prostor $R$-lineárních fcionálů $V\to R$ \\
\kw{reflexivní modul} & -- $V^{**}\cong V$
\end{tabular}
\begin{itemize}[nosep]
  \item[-] platí pro konečně generované moduly
  \item[-] platí pro řezy vekt.\ bundlů s kon.\ dim.\ fibrem
\end{itemize}

\begin{vysvetleni}[Moduly -- upřesnění]
Lineární závislost vyžaduje netriviální kombinaci (jinak by závislá byla každá $k$-tice). V modulu z ní neplyne vyjádřitelnost: z $r_1a_1=-r_2a_2$
nelze dělit $r_1$, není-li invertibilní (např.\ funkce s nulovými body). Standardně je \emph{báze} množina, která modul generuje \emph{a} je lineárně nezávislá
(rozklad $a=\sum a^ie_i$ je pak jednoznačný); bez nezávislosti jde jen o množinu generátorů. Konečně generovaný modul má \emph{konečnou množinu generátorů}
(nemusí mít bázi), volný modul s konečnou bází je $R^n$. Reflexivita obecně neplatí pro každý konečně generovaný modul (např.\ $\mathbb{Z}/2$ nad $\mathbb{Z}$ má
$V^*=0$), ale platí pro konečně generované \emph{projektivní} moduly -- a takové jsou podle Serreovy--Swanovy věty právě moduly řezů vektorových bandlů
(odtud druhý bod; např.\ $\mathcal{T}^*M\cong(\mathcal{T}M)^*$).
\end{vysvetleni}

%----------------------------------------------------------------------
\newpage
\strana{GMI-04b-7}
\begin{definice}[předsvazek $\mathcal{S}$ nad $M$ (též $\mathcal{S}_M$ či $\mathcal{S}M$) {[presheaf]}]
\begin{itemize}[nosep]
  \item topologický prostor $M$ \hfill (většinou top.\ varieta)
  \item přiřazení \quad $\mathcal{S}: U\subset M\to\mathcal{S}U$ \hfill ($U$ otevřené) \\
        kde $\mathcal{S}U$ je prostor $\mathcal{S}$-řezů nad $U$ \\
        typicky abel.\ gr., okruh, modul, algebra
  \item operace \kw{restrikce} $r^U_V$
  \[
    V\subset U \qquad r^U_V:\mathcal{S}U\to\mathcal{S}V \qquad \phi\mapsto\phi|_V \qquad\text{splňující}
  \]
  \begin{tabular}{@{}lll}
    $r^V_W\circ r^U_V = r^U_W$ & pro & $W\subset V\subset U\subset M$ otevř. \\
    $r^V_V = \mathrm{id}_V$ & pro & $V\subset M$ otevř.
  \end{tabular}
\end{itemize}
\end{definice}

\begin{vysvetleni}[Předsvazek]
Předsvazek je přiřazení \uv{objektů definovaných na otevřené množině $U$} (funkcí, polí, řezů \dots) každé otevřené $U$ spolu s operací zúžení
na menší množinu. Podmínky říkají jen, že zúžit dvakrát je totéž co zúžit najednou a že zúžení na tutéž množinu nic nemění
(v řeči kategorií: kontravariantní funktor z otevřených množin do grup, okruhů, \dots). Zatím nic neříká, že objekt je určen \uv{lokálními daty}.
\end{vysvetleni}

\begin{definice}[svazek $\mathcal{S}$ nad $M$ (též $\mathcal{S}_M$ či $\mathcal{S}M$) {[sheaf]}]
\kw{svazek} $\mathcal{S}_M$ je předsvazek $\mathcal{S}_M$ splňující
\begin{enumerate}[label=\arabic*)]
  \item $\phi_1,\phi_2\in\mathcal{S}U$ \quad $\forall\alpha$ \quad $\phi_1|_{U_\alpha} = \phi_2|_{U_\alpha}$ \quad $\Rightarrow$ \quad $\phi_1 = \phi_2$
  \item pro kolekci $\phi_\alpha\in\mathcal{S}U_\alpha$ splňující
  \[
    \forall\alpha,\beta \qquad \phi_\alpha|_{U_\alpha\cap U_\beta} = \phi_\beta|_{U_\alpha\cap U_\beta}
  \]
  existuje $\phi\in\mathcal{S}U$ tak, že $\phi|_{U_\alpha} = \phi_\alpha$
\end{enumerate}
kde $U = \bigcup_\alpha U_\alpha$, \quad $U_\alpha$ otevřené v $M$
\end{definice}

\begin{vysvetleni}[Axiomy svazku]
1) \emph{Lokalita}: řez je určen svými zúženími na libovolné otevřené pokrytí -- dva řezy, které se lokálně shodují, jsou stejné.
2) \emph{Slepování}: lokální řezy, které souhlasí na překryvech, lze slepit do řezu na sjednocení (díky 1) jednoznačně).
Svazek je tedy přesně ten typ objektu, který lze definovat \uv{lokálně a slepit} -- jako hladké funkce či vektorová pole, jejichž hladkost je lokální vlastnost.
Stejná logika jako u věty o rekonstrukci bandlu z přechodových zobrazení v GMI-04a.
\end{vysvetleni}

Pozn: pokud máme nějakou operaci mezi dvěma (před)svazky, předpokládáme konzistenci operace s restrikcí

\begin{priklad}[Př]
\begin{itemize}[nosep]
  \item[-] svazek $\mathcal{F}_M$ okruhů $C^\infty$ fcí na $M$ a jejích otevř.\ mn. \\
           máme řezy nad každou $U$ a i glob.\ rozšíření
  \item[-] svazek modulů $\mathcal{T}_M$ nad svazkem $\mathcal{F}_M$ \\
           tj.\ $\mathcal{T}U$ je $\mathcal{F}U$-modul \\
           nemusí existovat gl.\ nenul.\ řezy \quad (neparalelizov.\ $\mathbb{T}M$)
\end{itemize}
\end{priklad}

\begin{vysvetleni}[Pozor na \uv{globální rozšíření}]
Pro hladké funkce neplatí, že by \emph{každou} $f\in\mathcal{F}U$ šlo rozšířit na celé $M$: např.\ $f(x)=1/x$ na $U=(0,1)\subset\mathbb{R}$ hladké rozšíření
na $\mathbb{R}$ nemá. Platí ale (díky \uv{bump} funkcím, tj.\ rozkladu jednotky), že každou $f\in\mathcal{F}U$ lze po zúžení na menší okolí $V$ s $\overline V\subset U$
(tj.\ každý \emph{zárodek}, viz dále) rozšířit na globální hladkou funkci. To je velký rozdíl oproti holomorfním funkcím.
Totéž platí pro vektorová pole. \uv{Nenulové globální řezy} u $\mathcal{T}_M$: na $S^2$ neexistuje všude nenulové spojité vektorové pole (věta o \uv{chlupaté kouli}).
\end{vysvetleni}

%----------------------------------------------------------------------
\newpage
\strana{GMI-04b-8}
\begin{priklad}[Př: předsvazek $\mathcal{K}_M$ konst.\ fcí na $M$]
$\mathcal{K}U$ \quad konstantní fce na $U$ \\
\pozor{není svazek!}
\begin{itemize}[nosep]
  \item[] $M = M_1\cup M_2$ \qquad $M_i$ komponenty $M$
  \item[] řezy na $M_1, M_2$ nezaručují existenci řezu na $M$ \\
          $k_1\in\mathcal{K}M_1$, \ $k_2\in\mathcal{K}M_2$ \quad konzist.\ na $M_1\cap M_2 = \emptyset$ \\
          pro $k_1\neq k_2$ \pozn{indexy u $k_1, k_2$ nečitelné} neexistuje $k\in\mathcal{K}M$ tak, že
          $k|_{M_1} = k_1$ a $k|_{M_2} = k_2$
\end{itemize}
lze zavést svazek lokálně konst.\ fcí
\end{priklad}

\begin{vysvetleni}[Konstantní vs.\ lokálně konstantní funkce]
Selhává axiom slepování: na dvou disjunktních komponentách můžeme zvolit různé konstanty (souhlasí \uv{na prázdném překryvu}),
ale slepená funkce není konstantní. Funkce, které jsou konstantní jen na každé komponentě (lokálně konstantní), už svazek tvoří --
je to tzv.\ zesvazkování předsvazku $\mathcal{K}_M$ (konstantní svazek). Prostor jeho globálních řezů je $\mathbb{R}^{\#\text{komponent}}$.
\end{vysvetleni}

\begin{priklad}[Př: předsvazek $\mathcal{B}_M$ omezených fcí na $M$]
$\mathcal{B}U$ \quad soubor omezených hlad.\ fcí na $U$ \\
\pozor{není svazek!} \\
\quad omezenost na okolích $U_\alpha$ nezaručuje omezenost na jejich sjednocení -- pro nekonečně okolí
\end{priklad}

\begin{vysvetleni}[Příklad]
$f(x)=x$ na $M=\mathbb{R}$ je omezená na každém intervalu $U_n=(n-1,n+1)$ a zúžení na překryvy samozřejmě souhlasí, slepená funkce však na $\mathbb{R}$
omezená není. Omezenost je \emph{globální} vlastnost, kdežto svazky popisují vlastnosti, které jsou lokální (spojitost, hladkost, holomorfnost).
Pro konečná pokrytí by slepování fungovalo.
\end{vysvetleni}

\begin{priklad}[Př: svazek $\mathcal{H}$ holomorfních fcí na komplexní var.\ $M$]
$\mathcal{H}U$ \quad holomorfní fce na $U\subset M$ \\
nemusí existovat netriv.\ gl.\ řezy rozšiřující lokální řez, tj.\ pro $\phi_U\in\mathcal{H}U$ nemusí existovat
\oprava{$\phi\in\mathcal{H}M$}{$\phi\in M$} tak, že $\phi|_U = \phi_U$ \\
\quad např.\ $\sqrt z$, $\log z$, \dots \quad -- více větví
\end{priklad}

\begin{vysvetleni}[Holomorfní funkce a rozšíření]
Holomorfní funkce jsou \uv{tuhé}: podle věty o jednoznačnosti je holomorfní funkce na souvislé oblasti určena hodnotami na libovolně malé otevřené podmnožině.
Proto nelze použít bump funkce a lokální řez často nejde rozšířit: větev $\sqrt z$ či $\log z$ na oblasti kolem $z_0\neq0$ nelze rozšířit holomorfně
na $\mathbb{C}\setminus\{0\}$ (po oběhu počátku se změní větev), $1/z$ nelze rozšířit na $\mathbb{C}$. Na kompaktní souvislé komplexní varietě jsou dokonce
všechny globální holomorfní funkce konstantní.
\end{vysvetleni}

\begin{priklad}[Př: svazek $\mathcal{F}_M$ hladkých fcí na $M$]
$\mathcal{F}U$ \quad hladké fce na $U\subset M$ \\
pro každou $f\in\mathcal{F}U$ existuje globální rozšíření
\end{priklad}

\begin{vysvetleni}[Upřesnění]
Tvrzení je potřeba číst \uv{po zúžení na menší okolí} (resp.\ pro zárodky): $f\in\mathcal{F}U$ vynásobíme hladkou funkcí, která je $1$ na $V$ ($\overline V\subset U$)
a mimo kompaktní část $U$ je nulová, a dostaneme globální funkci shodnou s $f$ na $V$. Doslova pro každou $f$ to neplatí (viz $1/x$ na $(0,1)$, vysvětlivka výše).
\end{vysvetleni}

\begin{priklad}[Př: svazek $\mathcal{T}_M$ hladkých vekt.\ polí na $M$]
$\mathcal{T}U$ \quad hladká pole na $U\subset M$ \\
pro $U$ homeomorfní $\mathbb{R}^d$ je $\mathcal{T}U$ konečně gen.\ $\mathcal{F}U$-modul \\
ale $\mathcal{T}M$ nemusí být volný modul $=$ neparalelizovatelnost \\
$\mathcal{T}M$ je reflexivní (stačí lokálně konečně generovaný)
\end{priklad}

\begin{vysvetleni}[Vektorová pole jako modul]
Je-li $U$ obor mapy, je $\mathcal{T}U$ dokonce \emph{volný} modul s bází $\partial/\partial x^i$: každé pole je $a=a^i\partial/\partial x^i$ s jednoznačnými $a^i\in\mathcal{F}U$.
Globálně je $\mathcal{T}M$ volný právě tehdy, když existuje globální báze vektorových polí, tj.\ $\mathbb{T}M$ je paralelizovatelný ($S^1$, $S^3$, Lieovy grupy ano; $S^2$ ne).
I bez volnosti je $\mathcal{T}M$ konečně generovaný projektivní $\mathcal{F}M$-modul, a proto reflexivní: $(\mathcal{T}^*M)^*\cong\mathcal{T}M$ -- vektorová pole jsou
právě $\mathcal{F}M$-lineární funkcionály na 1-formách (srv.\ Krtouš, věta V2.3 o ultralokálních lineárních zobrazeních).
\end{vysvetleni}

%----------------------------------------------------------------------
\newpage
\strana{GMI-04b-9}
\begin{definice}[germ (zárodek) a stvol {[stalk]}]
mějme (před)svazek $\mathcal{S}_M$ a bod $x\in M$ \\
$\underset{x}{\sim}$ relace ekvivalence na řezech na oblastech obsah.\ $x$
\[
  \underset{\in\mathcal{S}U_1}{\phi_1}\ \underset{x}{\sim}\ \underset{\in\mathcal{S}U_2}{\phi_2}
  \quad\equiv\quad
  \exists\, U\subset M,\ x\in U,\ U\subset U_i:\quad \phi_1|_U = \phi_2|_U
\]
\hfill existuje okolí $U$, na kterém $\phi_1,\phi_2$ shodné

třída ekvivalence $\underset{x}{\sim}$ se nazývá \kw{germ} v bodě $x$ \\
soubor germů v bodě $x$ tvoří \kw{stvol} [stalk] $\mathcal{S}_xM$
\end{definice}

\begin{vysvetleni}[Zárodky]
Zárodek je \uv{nekonečně malý kousek} řezu kolem $x$: zapomeneme vše, co se děje mimo libovolně malé okolí $x$. Např.\ zárodky hladkých funkcí v $x$
nesou informaci o hodnotě a všech derivacích v $x$ (a ještě o něco víc), ale ne o funkci daleko od $x$. Stvol $\mathcal{F}_xM$ zárodků hladkých funkcí je okruh;
tečné vektory v $x$ lze definovat jako derivace tohoto okruhu. Stvol je algebraický objekt přiřazený jedinému bodu (\uv{vlákno} svazku nad $x$),
podobně jako vlákno bandlu.
\end{vysvetleni}

\begin{definice}[okruhované prostory {[ringed space]}]
$(M,\mathcal{R}_M)$ či jen $\mathcal{R}M$ je \kw{okruhovaný prostor} $\equiv$ \\
\quad $M$ \quad topologický prostor \\
\quad $\mathcal{R}_M$ \quad svazek okruhů

\kw{lokální okruhovaný prostor} $\mathcal{R}M$ $\equiv$ \\
\quad $\mathcal{R}_xM$ je lokální $=$ má jednoznačný maximální ideál
\end{definice}

\begin{vysvetleni}[Proč \uv{lokální} okruh]
Pro zárodky hladkých funkcí v $x$: je-li $f(x)\neq0$, je $f$ nenulová na nějakém okolí a $1/f$ tam existuje, tj.\ zárodek je invertibilní.
Neinvertibilní zárodky jsou tedy právě ty s $f(x)=0$ a tvoří ideál $\mathfrak{m}_x$ -- jediný maximální ideál stvolu $\mathcal{F}_xM$.
Bod $x$ je tak zakódován jako \uv{funkce, které v něm mizí}; navíc $T^*_xM\cong\mathfrak{m}_x/\mathfrak{m}_x^2$ (gradient zárodku mizejícího v $x$).
\end{vysvetleni}

umožňuje definovat pojem variety skrze zadání hladkých fcí

\begin{veta}[dif.\ varieta pomocí okruh.\ pr.]
$M$ je dif.\ var.\ (def.\ pomocí okruh.\ pr.) $\equiv$
\begin{itemize}[nosep]
  \item[] $(M,\mathcal{F}_M)$ je lokální okruh.\ prostor
  \item[] lokálně izomorfní s $(\mathbb{R}^n, C^\infty_{\mathbb{R}^n})$
  \item[] $M$ (s indukovanou top.) je Hausdorffův se spoč.\ bází top.
\end{itemize}
\end{veta}

umožňuje zobecnit pojem variety \\
\quad zadá se obecnější prostor fcí skrze zadání jiného lok.\ okruh.\ prostoru \\
\quad např.\ lokálně izomorfní s grassmannovským $(\mathbb{R}^{m,k}, C^\infty_{\mathbb{R}^{m,k}})$

\begin{vysvetleni}[Varieta jako lokálně okruhovaný prostor; supervariety]
\uv{Lokálně izomorfní} znamená: každý bod má okolí $U$ a homeomorfismus $U\cong V\subset\mathbb{R}^n$, který převádí svazek $\mathcal{F}_M|_U$ na svazek hladkých funkcí
na $V$ -- to je přesně mapa, jen formulovaná přes funkce místo souřadnic (souřadnice jsou $n$ vybraných funkcí). Hausdorffovost a spočetná báze vylučují
patologie (např.\ přímku se zdvojeným počátkem). Supervarieta: $\mathbb{R}^{m,k}$ (častěji psáno $\mathbb{R}^{m|k}$) je $\mathbb{R}^m$ se svazkem
$C^\infty(\mathbb{R}^m)\otimes\Lambda[\theta^1,\dots,\theta^k]$, kde $\theta^i$ jsou antikomutující (Grassmannovy) souřadnice, $\theta^i\theta^j=-\theta^j\theta^i$.
Body jsou stejné jako u $\mathbb{R}^m$, ale funkcí je \uv{víc} -- popisují fermionové stupně volnosti v supersymetrii.
\end{vysvetleni}

